\subsection{Noyau commun des \(3\) formulations de la dualité T}

Définissons le domaine stable
\[
 \mathscr D_0=\mathbb Z^2\times\mathbb R_{>0},
 \qquad
 E_{R_0}(m,w)=\frac{m^2}{R_0^2}+w^2R_0^2,
\]
et l'involution
\[
 \mathsf T_0(m,w,R_0)=(w,m,R_0^{-1}).
\]

\begin{theorem}[Équivalence du noyau ambiant reste--quotient]
\label{thm:cierre-vtm-tres-formulaciones-t}
Après prolongement des variables \(r,Q\) au domaine stable
\(\mathbb Z^2\), les trois formulations s'identifient au moyen de
\[
 I_F(r,Q,R)=(m,w,R_0)=(r,Q,R),
 \qquad I_{16\to02}=\operatorname{id}_{\mathscr D_0}.
\]
Ces applications entrelacent les involutions et les formes quadratiques :
\[
 I_F\mathsf T_F=\mathsf T_0I_F,
 \qquad
 E_R^F=E_{R_0}\circ I_F,
 \qquad
 \mathsf T_0^2=\operatorname{id},
 \qquad
 E_{R_0}=E_{R_0^{-1}}\circ\operatorname{swap}.
\]
\end{theorem}

\begin{proof}
Les deux applications sont des renommages des mêmes trois coordonnées. De plus,
\[
 \mathsf T_0^2(m,w,R_0)=(m,w,R_0)
\]
et
\[
 E_{R_0^{-1}}(w,m)
 =\frac{w^2}{R_0^{-2}}+m^2R_0^{-2}
 =w^2R_0^2+\frac{m^2}{R_0^2}
 =E_{R_0}(m,w).
\]
Les identités sont les mêmes après application de \(I_F\) ou de l'identité entre les
deux formulations.
\end{proof}

Dans \(\mathcal H_{\mathrm{lat}}=\ell^2(\mathbb Z^2)\), définissons
\[
 H(R_0)|m,w\rangle=E_{R_0}(m,w)|m,w\rangle,
 \qquad U_T|m,w\rangle=|w,m\rangle.
\]
Alors \(U_T\) est unitaire, \(U_T^2=I\), et l'identité scalaire précédente
équivaut, sur la base orthonormée, à
\[
 \boxed{U_TH(R_0)U_T^{-1}=H(R_0^{-1}).}
\]
C'est la deuxième réalisation exacte du même noyau, non une nouvelle
hypothèse physique.

\subsection{Distinction des 2 sections résiduelles par leurs insertions énergétiques et leurs classes d'équivalence}

La convention résiduelle positive utilise
\[
 r_9^+(n)=1+((n-1)\bmod9),
 \qquad Q_9^+(n)=\frac{n-r_9^+(n)}9,
\]
tandis que la convention résiduelle zéro insère la classe
\(\bar r_9(n)\in\{0,\ldots,8\}\). Le changement de section qui conserve l'
entier est
\[
 C(r,Q)=\bigl(\bar r,Q+\mathbf1_{\{r=9\}}\bigr).
\]

\begin{proposition}[Obstruction énergétique au changement de section]
\label{prop:cierre-vtm-obstruccion-secciones-nonadicas}
L'application \(C\) conserve \(n=9Q+r\), mais ne conserve pas en général la forme
\(E_R\) ni n'entrelace l'involution d'échange. Par conséquent, l'équivalence
du noyau ambiant n'implique pas l'équivalence des deux insertions
nonadiques.
\end{proposition}

\begin{proof}
Pour \(n=9\), la première convention donne \((r,Q)=(9,0)\), tandis que le changement avec retenue donne
\((\bar r,Q_0)=(0,1)\). Leurs énergies sont
\[
 E_R(9,0)=\frac{81}{R^2},
 \qquad E_R(0,1)=R^2,
\]
qui ne coïncident pas pour tout \(R>0\). L'insertion sans correction de
la retenue enverrait cet exemple sur \((0,0)\) et ne conserverait pas non plus l'entier.
Il faut fixer une section unique ou construire une forme covariante sous le
changement de section avant de promouvoir cette équivalence.
\end{proof}

\begin{definition}[Quotient de changement de section et énergie covariante]
\label{def:cierre-vtm-energia-covariante-seccion}
Soit
\[
 L_9:=\mathbb Z(9,-1)\subset\mathbb Z^2,
 \qquad
 S(x_1,x_2)=(x_2,x_1),
\]
et soit
\[
 q_R(x_1,x_2)=\frac{x_1^2}{R^2}+x_2^2R^2.
\]
Sur les quotients
\[
 \mathscr Q_9=(\mathbb Z^2/L_9)\times\mathbb R_{>0},
 \qquad
 \mathscr Q_9^\vee=(\mathbb Z^2/SL_9)\times\mathbb R_{>0},
\]
nous définissons
\[
 \mathcal E_R^{L_9}([x]_{L_9})
 :=\min_{\ell\in L_9}q_R(x+\ell)
\]
et
\[
 \mathsf T_{\mathrm{cov}}([x]_{L_9},R)
 :=([Sx]_{SL_9},R^{-1}).
\]
\end{definition}

\begin{theorem}[Équivalence covariante des sections nonadiques]
\label{thm:cierre-vtm-equivalencia-covariante-secciones}
L'énergie \(\mathcal E_R^{L_9}\) et la transformation
\(\mathsf T_{\mathrm{cov}}\) sont bien définies. Les conventions
\((9,0)\) et \((0,1)\), et en général les sections \(1,\ldots,9\) et
\(0,\ldots,8\) avec la retenue correspondante, représentent la même classe
de \(\mathbb Z^2/L_9\). De plus,
\[
 \boxed{
 \mathcal E_{R^{-1}}^{SL_9}([Sx]_{SL_9})
 =\mathcal E_R^{L_9}([x]_{L_9}).}
\]
L'échange inverse, défini sur \(\mathscr Q_9^\vee\), restitue la
classe initiale ; les deux cartes quotientées sont donc T-duales.
\end{theorem}

\begin{proof}
Si \(x'=x+\ell_0\), la translation
\(\ell\mapsto\ell+\ell_0\) permute \(L_9\), donc le minimum ne dépend pas du
représentant. Il existe parce que \(q_R(x+k(9,-1))\) est une fonction quadratique
coercive de \(k\in\mathbb Z\). La différence
\((9,0)-(0,1)=(9,-1)\) appartient à \(L_9\), et il en va de même à chaque changement
de section avec retenue. Enfin,
\[
\begin{aligned}
 \mathcal E_{R^{-1}}^{SL_9}([Sx])
 &=\min_{\ell\in L_9}q_{R^{-1}}(Sx+S\ell)\\
 &=\min_{\ell\in L_9}q_R(x+\ell)
 =\mathcal E_R^{L_9}([x]).
\end{aligned}
\]
Comme \(S^2=I\), le second échange retrouve le quotient initial et le
rayon \(R\).
\end{proof}

\paragraph{Domaine d'équivalence du noyau dual et non-équivalence des insertions énergétiques.}
L'équivalence des noyaux sur \(\mathscr D_0\) est exacte. Les deux
insertions ne sont pas équivalentes pour l'énergie non corrigée ; la clôture
covariante est le quotient et l'énergie du théorème précédent.

\subsection{Normalisation exacte de la réalisation de corde}

Fixons \(\alpha'>0\). Dans le domaine
\(\mathscr D_{\alpha'}=\mathbb Z^2\times\mathbb R_{>0}\), soit
\[
 \mathsf T_{\alpha'}(m,w,R)=\left(w,m,\frac{\alpha'}R\right),
 \qquad
 \mathcal N_{\alpha'}(m,w,R)
 =\left(m,w,\frac R{\sqrt{\alpha'}}\right).
\]

\begin{theorem}[Conjugaison par l'échelle de corde]
\label{thm:cierre-vtm-normalizacion-alpha-t}
La normalisation \(\mathcal N_{\alpha'}\) est une conjugaison exacte :
\[
 \boxed{\mathcal N_{\alpha'}\mathsf T_{\alpha'}
 =\mathsf T_0\mathcal N_{\alpha'}.}
\]
Pour
\[
 E_{\mathrm{str}}(m,w,R)
 =\frac{m^2}{R^2}+\frac{w^2R^2}{\alpha'^2},
\]
on a
\[
 E_{\mathrm{str}}
 =\alpha'^{-1}E_0\circ\mathcal N_{\alpha'}.
\]
De plus, si
\[
 p_L=\frac mR+\frac{wR}{\alpha'},
 \qquad p_R=\frac mR-\frac{wR}{\alpha'},
\]
alors \(p_L\) reste fixe et \(p_R\mapsto-p_R\).
\end{theorem}

\begin{proof}
On a
\[
 \mathcal N_{\alpha'}\mathsf T_{\alpha'}(m,w,R)
 =\left(w,m,\frac{\sqrt{\alpha'}}R\right)
 =\mathsf T_0\mathcal N_{\alpha'}(m,w,R).
\]
L'identité des énergies résulte de
\[
 E_0\!\left(m,w,\frac R{\sqrt{\alpha'}}\right)
 =\frac{\alpha'm^2}{R^2}+\frac{w^2R^2}{\alpha'}
 =\alpha' E_{\mathrm{str}}(m,w,R).
\]
La substitution de
\((m,w,R)\mapsto(w,m,\alpha'/R)\) dans \((p_L,p_R)\) donne respectivement
\(p_L\) et \(-p_R\).
\end{proof}

La conjugaison et l'invariance du spectre à données oscillatoires fixées sont
exactes. L'identification de \(m,w,R,\alpha'\) à l'impulsion, à l'enroulement,
au rayon physique et à la tension inverse est une interface physique munie de ses propres prémisses.

\subsection{Niveau, échange linéaire de masses et modularité : \(3\) contrôles de non-fusion}

La condition active de corde est
\[
 N_L-N_R=a_L-a_R-mw.
\]
La formule schématique \(N_L-N_R=nw\) ne coïncide avec elle
qu'après déclaration de \(a_L=a_R\) et de \(n=-m\), ou d'une inversion équivalente de
l'orientation des feuillets. Sans ce dictionnaire, les deux expressions ne sont pas la
même formulation.

La symétrie linéaire de l'action est typée dans l'espace étendu des
coefficients et des modes :
\[
 L(A,C;n,w)=nA+wC,
 \qquad
 \sigma(A,C;n,w)=(C,A;w,n).
\]
Alors
\[
 L\circ\sigma=wC+nA=L.
\]
C'est une identité exacte d'échange simultané des paramètres et des
coordonnées. Elle n'équivaut pas à l'invariance quadratique de \(H(R)\) sur un
secteur à paramètres physiques fixés ; une telle promotion requiert un
entrelaceur supplémentaire.

Enfin, sur l'axe imaginaire du demi-plan supérieur, définissons
\[
 \iota(R_0)=iR_0^2,
 \qquad S(\tau)=-\frac1\tau.
\]
Alors
\[
 \boxed{S\circ\iota(R_0)=\iota(R_0^{-1}).}
\]
Comme \(J(S\tau)=J(\tau)\), on obtient
\(J(iR_0^2)=J(iR_0^{-2})\). C'est une semi-conjugaison exacte du paramètre
de rayon avec la transformation modulaire \(S\), démontrée avec la structure
modulaire de départ explicite. Elle n'identifie pas à elle seule la dualité de l'
espace cible, l'échange des charges et la construction de
\(V^\natural\).
